\section{LITERATURE REVIEW}
\label{sec:literature}

This section narrows from the format itself to the systems the proposed pipeline
is positioned against. We first make explicit the mismatch between the
presentational structure of PDF and the semantic structure machine translation
requires (Section~\ref{subsec:lit-pdf}); we then review the model families the
pipeline builds on (Section~\ref{subsec:lit-models}); next we distinguish
document-conversion systems from document-translation systems
(Sections~\ref{subsec:lit-conversion} and~\ref{subsec:lit-systems}); and finally
we state the gap that motivates the hypothesis of Section~\ref{subsec:hypothesis}
and the design decisions of Section~\ref{sec:methodology}
(Section~\ref{subsec:lit-gap}).

\subsection{PDF Structure and Layout Preservation}
\label{subsec:lit-pdf}

\paragraph{Format characteristics.} Under ISO~32000-2, page content is a stream
of graphics operators that place glyphs, strokes and images into a fixed
coordinate system~\citep{pdf-spec}. This is appropriate for display and print,
but it does not guarantee that semantic units---headings, paragraphs, table
cells, reading order---can be recovered directly. With respect to content
accessibility, PDF documents fall into three groups. \emph{Born-digital} files
contain text objects with font and position information, and can be extracted
directly, although the order in which objects appear in the content stream need
not be the reading order. \emph{Scanned} files contain page bitmaps and no
trustworthy text layer, and therefore require OCR. \emph{Mixed} documents
contain born-digital text, scanned pages and text-bearing images at the same
time, so the extraction strategy must be decided per page rather than per
document.

Consequently, raw extraction is not sufficient input for machine translation.
An intermediate representation must record content, element type, bounding box,
page number and reading-order relations simultaneously; it is the bridge
between the semantic domain of translation and the geometric domain of
typesetting.

\paragraph{Element types and per-type policy.} Document layout analysis (DLA)
assigns positions and semantic labels to page regions. Datasets such as
DocLayNet provide bounding-box labels for many element types across diverse
document domains and thus support both training and
evaluation~\citep{doclaynet}, and surveys report that DLA quality directly
determines downstream information extraction, reading-order recovery and format
conversion~\citep{doclayout-survey}. For translation, each element group
requires a distinct policy: flowing text must be grouped into units with enough
context; formulas must preserve their symbols, with only the accompanying prose
translated; figures are copied unchanged while their captions are translated;
and table content must remain bound to the correct row and column, since
translating the flattened string of a table destroys inter-cell relations.

\paragraph{Text overflow.} The length and typesetting behaviour of a
translation differ from those of the source by an amount that depends on the
language pair, the content, the terminology, the font and the line-breaking
rules; a character count alone therefore does not predict the display space
required. When target text is placed into the original box, three failures
appear: content spills outside the box, it overlaps a neighbouring element, or
the font is reduced until legibility suffers. This creates a trade-off between
geometric fidelity and readability. Preserving coordinates and box sizes
exactly keeps the target page close to the source but maximises overflow risk;
re-flowing the page lets text adapt but moves figures, tables and page breaks.
A layout-preserving system must therefore combine several signals---box size,
source typography, surrounding free space, and the possibility of collision
with other elements---and its evaluation must consider layout similarity,
overlap count and readability rather than merely whether an output file was
produced.

\paragraph{Appearance beyond geometry.} Placing target text at the right
coordinates is necessary but not sufficient, because the replaced text also
carried a visual identity. Two properties are routinely lost. First, the
\emph{background}: writing a translation over a source page requires masking
what was there, and masking with a global white rectangle is correct only for
white pages---on coloured slides, textbook sidebars, shaded table rows or
scanned paper the mask is itself a visible defect. Second, the \emph{text
colour}: headings, hyperlinks, emphasised terms and slide text are frequently
non-black, and a reconstruction that repaints everything in black silently
discards a layer of document semantics. Neither property can be read reliably
from the content stream in the general case---scanned pages have no text objects
at all, and born-digital pages express colour through nested graphics
states~\citep{pdf-spec}---so a pipeline that already rasterises every page has
the option, largely unexploited in existing systems, of \emph{estimating}
appearance from the rendered pixels.

\subsection{Foundation Models for Document Analysis}
\label{subsec:lit-models}

A layout-preserving translation pipeline must coordinate a document vision
model, an OCR model, a translation model and a service-orchestration layer.
These components do not operate independently: a recognition error at the front
of the pipeline propagates into the translation and finally manifests as a
typesetting defect. Model selection must therefore weigh accuracy, language
coverage, inference cost, and---critically---whether the output format is usable
by the following stage.

\paragraph{Layout detection.} A layout detector consumes a page image and
returns regions with semantic labels and confidences. DocLayout-YOLO balances
accuracy against speed using diverse synthetic pre-training data and a
global-to-local perception mechanism, which suits document pages where a
one-line heading and a near-full-page table can co-occur~\citep{doclayout-yolo}.
Surya provides layout detection, reading-order estimation and table recognition
inside a single toolkit~\citep{surya}. For our pipeline the important output is
not only the class label but a stable bounding box and a reading order that can
be joined reliably with OCR results.

\paragraph{OCR for scanned PDFs.} A document OCR workflow typically detects
text lines or regions, recognises character sequences, and maps the result back
into page coordinates. Unlike OCR intended only for search, reconstructing a
PDF additionally requires line- or word-level boxes, recognition confidences
and the relation between text and layout regions. Surya supports multilingual
recognition, line detection, reading order, layout, tables and formulas in one
toolkit~\citep{surya}; PaddleOCR contributes multilingual recognition models and
structured document components, which are convenient for table-cell
information~\citep{paddleocr}. No single model is accurate on every document
type: resolution, skew, background noise, typeface and inline formulas all
degrade results. A pipeline must therefore normalise coordinate systems, retain
confidences and keep stable identifiers per element so that errors remain
diagnosable across stages.

\paragraph{Evaluation metrics.} For translation, lexical metrics such as
BLEU~\citep{bleu} and chrF~\citep{chrf} measure surface overlap, whereas neural
metrics such as COMET~\citep{comet} encode source, hypothesis and reference with
a multilingual encoder and regress a continuous adequacy score. The WMT metrics
shared task reports that neural metrics correlate better with human judgement
and are more robust than lexical ones~\citep{wmt-metrics-shared-task}, which is
why we treat COMET as primary and chrF++ as a diagnostic baseline. For document
parsing, OmniDocBench provides 1{,}651 richly annotated pages together with a
composite score built from text edit distance, table TEDS and formula
CDM~\citep{omnidocbench}; CDM itself scores a formula by rendering it and
matching detected symbols, which makes it invariant to notation differences
that penalise string-based metrics unfairly~\citep{cdm}. Most of the string
metrics used below reduce to the Levenshtein edit
distance~\citep{levenshtein}.

\subsection{Document Parsing and Conversion Systems}
\label{subsec:lit-conversion}

Several recent systems address the front half of the problem: turning PDFs into
machine-readable structured representations. Docling converts PDFs into a
document model that can be exported to Markdown or JSON, combining layout
analysis and table-structure recognition in a resource-conscious
package~\citep{docling}. Marker similarly targets Markdown, JSON, HTML and chunk
outputs, using Surya-based OCR and layout models with optional LLM correction for
tables, inline mathematics and block-level repair~\citep{marker}. MinerU focuses
on high-precision document content extraction across diverse layouts, including
tables, formulas and scanned or garbled PDFs~\citep{mineru}. olmOCR takes a
large-scale VLM route, linearising PDFs into clean text while preserving reading
order and structured content for language-model data pipelines~\citep{olmocr}.

These systems are important comparators for the \emph{parsing} side of this
work, and they motivate treating structure analysis as a first-class component
rather than a utility function. They are not, however, direct baselines for the
complete task studied here. Their primary outputs are structured text
representations, not translated PDFs composited back onto the source page; they
optimise readability, retrieval or data preparation rather than preservation of
the original page geometry after translation. A fair end-to-end comparison would
therefore require an additional reconstruction layer on top of each converter,
plus an annotated PDF benchmark that scores both translation adequacy and final
visual fidelity. In the absence of such a benchmark, the present paper compares
the parsing stage to public ground truth and treats direct system-level
comparison as future work rather than evidence of superiority.

\subsection{Related Document Translation Systems}
\label{subsec:lit-systems}

Document translation systems can be distinguished by the intermediate
representation they adopt: coordinate-based representations favour geometric
fidelity, structured-text representations favour context and editability, and
re-typesetting representations rebuild the document from recognised components.
The comparison below concerns published architecture and scope, and does not
treat feature count as evidence of output quality. The systems are discussed
qualitatively because their public implementations do not expose the same input
scope and output contract on scanned, mixed and born-digital PDFs.

\paragraph{Layout preservation by text overlay.} PDFMathTranslate~\citep{pdf2zh}
implements layout detection, content extraction, translation and document
reconstruction. A layout model separates text from regions that must be
preserved, and the text is passed to an intermediate layer supporting several
translation services; target content is written back using the geometry of the
source, so figures, formulas and page structure are disturbed far less than in
a plain-text conversion. The important contribution is the separation of the
layout pipeline from the translation service, which allows the language model to
be replaced without redesigning the PDF parser. Output quality nevertheless
remains bounded by layout-detection accuracy and by the ability to fit
new-length text into old regions, and the published version identifies
suboptimal OCR support as a limitation, which restricts applicability to
scanned PDFs~\citep{pdf2zh}.

\paragraph{Segmentation and terminology consistency.}
docutranslate~\citep{docutranslate} approaches the task from the content side:
it supports many formats, splits content into segments compatible with model
input limits, executes translation asynchronously and supplies a glossary.
Terminology constraints matter particularly for long documents, where the same
concept must be rendered consistently across segments that are never processed
in the same call; their effect has been studied in machine
translation~\citep{terminology-mt}, and recent work on document-level
translation reports that fusing several knowledge sources improves contextual
consistency~\citep{doc-mt-multi-knowledge}. For PDF input, docutranslate uses a
document parser to recognise tables, formulas and code, then converts content to
Markdown before translation. This normalisation is convenient for semantic
processing and multi-format export, but the project documentation states that
the original PDF layout is lost in the conversion~\citep{docutranslate}.

\paragraph{Re-typesetting from OCR.} RetainPDF~\citep{retain-pdf} represents the
OCR-and-rebuild direction, covering content recognition, translation, diagnosis
and PDF export, with an input scope that includes image-only and scanned
documents. Compared with exploiting the existing text layer, this yields a more
uniform representation of born-digital and scanned PDFs and allows control over
fonts, tables and formulas during typesetting. In exchange, rebuilding from OCR
introduces a chain of accumulated error: mis-detected regions produce wrong
reading order, mis-recognised characters change the translation input, and
geometric drift affects placement at export time. As the system is described
mainly through its source repository rather than a peer-reviewed empirical
study, its quality claims require verification on a common dataset and a common
metric set before direct comparison.


\begin{table}[H]
    \centering
    \caption{Qualitative comparison of related document translation approaches.}
    \label{tab:related-work-comparison}
    \footnotesize
    % Fixed 2.6cm + 3x4.0cm columns were set for the 6.75in two-column measure
    % and now overrun the text block; X columns divide whatever measure is in
    % force, so the table fits by construction.
    \begin{tabularx}{\textwidth}{@{}
        >{\raggedright\arraybackslash}p{2.3cm}
        >{\raggedright\arraybackslash}X
        >{\raggedright\arraybackslash}X
        >{\raggedright\arraybackslash}X@{}}
        \toprule
        \textbf{Criterion} & \textbf{PDFMathTranslate} & \textbf{docutranslate} & \textbf{RetainPDF} \\
        \midrule
        Representation
        & Bound to the source PDF
        & Markdown or format-specific structure
        & Re-typeset PDF \\
        Primary objective
        & Preserve formulas and page layout
        & Translation context, terminology, many formats
        & Support image/scanned PDF and re-export \\
        Stated limitation
        & OCR support suboptimal in the published version
        & PDF-to-Markdown conversion loses the original layout
        & Error may accumulate through OCR and rebuilding \\
        Published evidence
        & Peer-reviewed system paper
        & Documentation and source repository
        & Documentation and source repository \\
        \bottomrule
    \end{tabularx}
\end{table}

\subsection{Synthesis and Research Gap}
\label{subsec:lit-gap}

The systems above are not mutually exclusive; they optimise different
objectives. PDFMathTranslate prioritises the geometric layout of born-digital
PDFs, docutranslate prioritises a content representation convenient for
document-level translation, RetainPDF extends coverage to scanned documents
through OCR and reconstruction, and document-conversion systems such as
Docling, Marker, MinerU and olmOCR prioritise high-quality structured extraction
for downstream applications. Table~\ref{tab:related-work-comparison} summarises
the differences drawn from published papers and project documentation.

The gap should therefore not be phrased as a single missing feature---several
tools already support scanned PDFs, structured extraction or layout-aware export
to varying degrees. What is still under-measured is the \emph{joint} problem of
translation, geometry and appearance: the translated content must remain linked
to page coordinates, and the reconstruction stage must be evaluated as part of
the task rather than treated as a cosmetic post-process. Four requirements are
therefore not addressed simultaneously by the surveyed systems:

\begin{itemize}[leftmargin=1.2em]
  \item \textbf{Unified input handling.} Born-digital, scanned and mixed PDFs
  should follow one processing path while the provenance of each extracted
  result remains traceable.
  \item \textbf{Joint semantics and geometry.} Segmentation and glossaries
  improve translation consistency, but the translated content must be mapped
  back onto individual geometric elements rather than existing only as text.
  \item \textbf{Page-aware overflow handling.} Uniform shrinking degrades
  readability, while unconstrained overflow produces overlap. The scaling
  decision must depend on whether a collision with a neighbouring element would
  actually occur---which the translator, having no access to the page, cannot
  determine.
  \item \textbf{Appearance fidelity.} A translated block must inherit the
  background it sits on and the colour it is drawn in. Neither is available from
  the content stream for scanned input, so both have to be estimated from the
  rendered page if the output is to remain faithful on anything other than
  black-on-white documents.
\end{itemize}

Two consequences follow, and they structure the rest of the paper. First, because
the surveyed tools report evidence in different forms, no conclusion about which
stage limits output quality can be drawn from feature descriptions alone. That
question must first be approached by measuring the stages apart, which is the
protocol of Section~\ref{subsec:method-eval} and the substance of hypothesis~H1.
Second, because layout preservation is a property of the page rather than of the
text, the requirements above are discharged at reconstruction time, which is
hypothesis~H2 and the substance of Section~\ref{subsec:method-render}.
